\chapter{Panel data: heterogeneity, timing and structural moments}
\label{r3:panels}
\section{Repeated measurements do not specify the identification strategy}
A panel of branch deposits can support several different questions.
One can model a homogeneous response to price after removing stable
branch differences, identify an adoption-cohort effect from untreated
trends, fit an interactive-factor counterfactual, or estimate a dynamic
equation using lagged instruments. All use a rectangular data set;
their counterfactual assumptions are different.

Let $i=1,\ldots,N$ index branches and $t=1,\ldots,T$ calendar periods.
Write $Y_{it}(\bar a_t)$ for the outcome under an action history.
The shorter $Y_{it}(a_t)$ assumes away carryover. Let $G_i$ be the
first adoption date only when an absorbing intervention makes that
description meaningful. A price that rises and falls, a repeated
collections action and a permanent branch rollout do not share the
same treatment history merely because each is coded as a binary column.
\citet[Sections 2--4]{panelreview2024} emphasizes assignment pattern,
potential-outcome dynamics and the relative sizes of $N,T$ as distinct
organizing dimensions.

Banking measurements also change composition. Branch closure can remove
a low-performing unit, household transfers can move balances between
branches, and a merger can change the entity represented by an identifier.
An unbalanced panel is not automatically missing at random. The outcome
and unit definitions of Chapter~\ref{r3:graphs} remain part of the
panel estimand.

\section{Fixed effects remove a component, not all confounding}
\label{r3:panel-fe}
Consider the static linear structural specification
\begin{equation}
 Y_{it}=\alpha_i+\lambda_t+\beta A_{it}+X_{it}'\gamma+u_{it}.
 \label{r3:static-panel}
\end{equation}
The interpretation of $\beta$ as a common causal response requires
linearity, absence of omitted carryover and an invariant outcome
equation. A sufficient mean restriction for the within estimator is
strict exogeneity,
\[
 \E[u_{it}\mid \alpha_i,A_{i1:T},X_{i1:T}]=0
 \quad\text{for every }t,
\]
with common time components treated appropriately. This permits
$\alpha_i$ to correlate with the entire action/covariate history.
It rules out current shocks predicting future actions after conditioning
on the specified regressors. For example, an unexpected deposit outflow
that causes next month's retention offer threatens strict exogeneity,
even if the offer cannot affect past outcomes.

For a balanced panel, two-way demeaning gives
$\ddot A_{it}=A_{it}-\bar A_{i\cdot}-\bar A_{\cdot t}+\bar A$,
and analogously for $Y,X,u$. The stable branch effect and common time
effect vanish. Frisch--Waugh residualization by $\ddot X$ then gives
\[
 \widehat\beta-\beta
 =\frac{\sum_{it}\widetilde A_{it}\ddot u_{it}}
        {\sum_{it}\widetilde A_{it}^2}.
\]
Identification needs the numerator's mean restriction and nonzero
residual variation; the data transformation supplies neither.
For unbalanced panels, projecting on actual entity and time design
matrices replaces this simple double-demeaning formula. Treating
missing cells as zero changes that projection.

\subsection{Random and correlated effects change the restrictions}
Random-effects GLS uses between-unit as well as within-unit variation.
In the elementary model with $\alpha_i$ independent of regressors,
$\Var(\alpha_i)=\sigma_\alpha^2$ and independent homoskedastic
$u_{it}$, the within-unit covariance is
$\Omega_i=\sigma_u^2I_T+\sigma_\alpha^2\mathbf1\mathbf1'$.
Its eigenvalue in the constant direction is
$\sigma_u^2+T\sigma_\alpha^2$, and in the remaining directions
it is $\sigma_u^2$. Hence, up to a common scale, GLS quasi-demeans
each variable as $x_{it}-\vartheta\bar x_i$, with
\[
 \vartheta=1-\sqrt{\frac{\sigma_u^2}{\sigma_u^2+T\sigma_\alpha^2}}.
\]
This project covariance derivation explains its efficiency gain under
the stronger mean model and its exposure to confounding between units.

A correlated-effects specification instead posits, for example,
$\E[\alpha_i\mid A_{i1:T},X_{i1:T}]
=c+\delta\bar A_i+\bar X_i'\pi$ and adds those means to the
regression. It permits a particular dependence on histories while
restricting its functional form. It does not make arbitrary adaptive
assignment ignorable. Nonlinear models also introduce an incidental
parameter issue when one estimates many unit intercepts with small
$T$; removing linear means is not a general nonlinear solution
\citep[Sections 2 and 3]{panelreview2024}.

\section{Dynamics turn lagged outcomes into endogenous regressors}
\label{r3:dynamic-panel}
Consider
\begin{equation}
 Y_{it}=\rho Y_{i,t-1}+\beta A_{it}+\alpha_i+\lambda_t+u_{it}.
 \label{r3:dynamic-panel-model}
\end{equation}
Even if $u_{it}$ is a martingale difference relative to past outcomes,
the within transform of $Y_{i,t-1}$ contains earlier innovations
also present in the within transform of the error. Differencing gives
\[
 \Delta Y_{it}=\rho\Delta Y_{i,t-1}+\beta\Delta A_{it}
                        +\Delta\lambda_t+\Delta u_{it}.
\]
In the simplest case with no action, $\alpha_i=0$, $Y_{i0}=0$
and independent unit-variance $u_{i1},u_{i2}$, the two-period
regression uses $\Delta Y_{i1}=u_{i1}$ and
$\Delta Y_{i2}=(\rho-1)u_{i1}+u_{i2}$. Its population coefficient
is $\rho-1$, not $\rho$. More branches do not remove this short-panel
bias. This exact construction illustrates the problem without claiming
that its numerical bias formula applies at arbitrary $T$.

\subsection{The valid lag is determined by a moment restriction}
If innovations are serially uncorrelated and initial observations are
orthogonal to future innovations, then for appropriate $s\le t-2$
\begin{equation}
 \E[Y_{is}(u_{it}-u_{i,t-1})]=0.
 \label{r3:ab-moments}
\end{equation}
By contrast, $Y_{i,t-1}$ contains $u_{i,t-1}$ and is not a valid
instrument for the differenced equation. For action regressors, strictly
exogenous $A_{is}$ can supply all date moments. Predetermined actions
orthogonal to current/future innovations can supply $s\le t-1$;
actions correlated with current innovations require earlier valid
lags, typically $s\le t-2$ under the corresponding sequential
restriction. These are assumptions about the assignment/error clock,
not a software convention about lag length.

Stack valid instruments in $Z_i$ and differenced structural residuals
in $\Delta u_i(\theta)$, with $\theta=(\rho,\beta,\ldots)$.
The GMM criterion is
\[
 \widehat\theta=\arg\min_\theta
 \left\{N^{-1}\sum_i Z_i'\Delta u_i(\theta)\right\}'\widehat W
 \left\{N^{-1}\sum_i Z_i'\Delta u_i(\theta)\right\}.
\]
This is the moment-estimation framework of
\citet[Section 2, equations (2)--(6)]{arellano1991}. Its sampling
unit is an independent individual time series, with fixed $T$ in the
main argument. Differenced white-noise innovations have first-order
covariance $-\sigma_u^2$ and zero covariance at lag two; thus first-order
correlation in differenced residuals is expected, whereas a second-order
correlation diagnostic probes a key validity restriction. The paper's
Section 3 derives the test and its estimation corrections. Failure to
reject does not establish the exclusion of a bank's omitted policy
response or measurement process.

Persistent balances can make lagged levels weak instruments for
changes. A large menu of lags can also create many moments relative to
the number of independent branches. The weak-IV discussion then
matters inside the panel estimator. Extra level-equation moments, for
example $\E[\Delta Y_{i,t-1}(\alpha_i+u_{it})]=0$, require additional
initial-condition/mean-stationarity restrictions; they do not follow
from Equation~\eqref{r3:ab-moments}. A product option that adds such
moments must expose those extra restrictions explicitly. Section~\ref{s5:dynamic-panel} develops the additional initial-law
argument and GMM covariance, and distinguishes the implemented first-order
sandwich from an unimplemented finite-sample correction.

\section{Difference-in-differences identifies an untreated evolution}
\label{r3:did-identification}
In a two-period rollout let $D_i$ indicate treatment in the second
period and assume no one is treated in the first. For this subsection
write $\Delta Y=Y_1-Y_0$. Under no anticipation,
\[
 \operatorname{ATT}=\E[Y_1(1)-Y_1(0)\mid D=1].
\]
Conditional parallel trends imposes
\begin{equation}
 \E[Y_1(0)-Y_0(0)\mid D=1,X]
 =\E[Y_1(0)-Y_0(0)\mid D=0,X]=m_0(X).
 \label{r3:parallel-trends}
\end{equation}
It does not require equality of untreated levels. Subtract the
counterfactual untreated change within the treated population to get
\begin{equation}
 \operatorname{ATT}=\E[\Delta Y-m_0(X)\mid D=1].
 \label{r3:did-functional}
\end{equation}
This is a different identification route from
$Y(a)\perp A\mid X$ or sequential randomization. Neither implies
the other in general. Differencing can remove a stable latent level,
but not an unmeasured shock whose differential trend also drives
adoption. Adding a post-rollout balance to $X$ can additionally change
the causal conditioning set.

Let $p(X)=P(D=1\mid X)$ and $p_D=P(D=1)>0$. With overlap for
the control group on the treated covariate support,
\begin{equation}
 \operatorname{ATT}
 =\frac1{p_D}\E\!\left[
 \left\{D-(1-D)\frac{p(X)}{1-p(X)}\right\}
                       \{\Delta Y-m_0(X)\}\right].
 \label{r3:did-dr}
\end{equation}
Conditioning on $X$ proves the formula. If a candidate $m$ equals
$m_0$, its control residual has zero mean for any candidate propensity.
If a candidate propensity equals $p$, its weighting balances the
$X$ terms, cancelling an arbitrary candidate $m$. More precisely,
using candidates $\widetilde p,\widetilde m$ and the true normalizer,
the bias is
\[
 \frac1{p_D}\E\!\left[
 \frac{\widetilde p-p}{1-\widetilde p}
                         (\widetilde m-m_0)\right].
\]
This reconstructs the mechanism of the outcome, weighting and
DR representations in \citet[Section 2.4 and Theorem 1]{callaway2021}.
Estimated normalizers and aggregation weights also enter its sampling
influence function. A generic ATE standard-error routine does not
automatically supply ATT or event-study inference.

\section{Staggered adoption requires explicit comparison groups}
\label{r3:staggered}
For absorbing adoption date $G_i=g$, define
$\operatorname{ATT}(g,t)=\E[Y_{it}(g)-Y_{it}(\infty)\mid G_i=g]$.
With no anticipation and valid conditional untreated trends, use
baseline $g-1$ and a comparison group $C_{g,t}$ that is untreated
through $t$. A never-treated group and a not-yet-treated group impose
different trend restrictions and offer different support. If effects
can begin $\delta$ periods before recorded adoption, use baseline
$g-\delta-1$ and controls unexposed to anticipation through the
comparison horizon. These are the separate Assumptions 3--6 and
identification routes in \citet{callaway2021}.

Apply Equation~\eqref{r3:did-functional} to each eligible group/time
comparison, with its own covariate overlap, outcome difference and
control law. Then state the aggregation. An event-time curve
$\sum_g w_{g,k}\operatorname{ATT}(g,g+k)$ can change composition
with $k$ because later cohorts lack long follow-up. A balanced-cohort
curve uses common cohorts and different information. The product must
return the cohort weights and supported horizons, not just a line
labelled ``dynamic effect.''

\subsection{A negative coefficient with positive effects everywhere}
Two-way fixed-effects regression projects treatment on unit/time
indicators and uses residual $\widetilde D_{it}$. If untreated outcomes
have that additive form and treatment effects are $\tau_{it}$, its
population coefficient in a fixed design is
\begin{equation}
 \beta_{\rm TWFE}
 =\sum_{it:D_{it}=1}\omega_{it}\tau_{it},\qquad
 \omega_{it}=\frac{\widetilde D_{it}}
 {\sum_{js:D_{js}=1}\widetilde D_{js}}.
 \label{r3:twfe-weights}
\end{equation}
The weights sum to one, but need not be positive. The identity follows
from projection orthogonality and
$\sum\widetilde D D=\sum\widetilde D^2$; it is a direct
decomposition, not a causal justification for the weights.

Take three equally weighted branches and four periods, with treatment
histories $(0,1,1,1)$, $(0,0,0,1)$ and $(0,0,0,0)$. Set all
untreated outcomes to zero. The four treated-cell effects, ordered by
early periods 2, 3, 4 and late period 4, are $(1,1,20,1)$.
The corresponding TWFE weights are $(0.3,0.3,-0.1,0.5)$, giving
$\beta_{\rm TWFE}=-0.9$ although every effect is positive.
The equally weighted treated-cell ATT is $23/4=5.75$.
Group/time differences against the never-treated branch recover the
four effects exactly in this noiseless construction. The example
isolates target weighting; it is not a small-sample noise effect.
The heterogeneous-effects issue and alternative constructions are
developed in \citet{callaway2021,panelreview2024}.

If actions can reverse, first-adoption coding discards part of the
intervention path. One may instead compare a switching group's
realized path after its first change with maintaining its initial
treatment, using groups with the same initial treatment that have not
yet switched. No anticipation and parallel evolution of these
status-quo outcomes are new restrictions. The intertemporal estimators
of \citet[Assumptions 1--2 and Section 3]{dechaisemartin2024} develop
that route and make normalization across treatment increments explicit.
It does not identify arbitrary histories without support. For repeated
feedback-driven banking actions, the longitudinal g-method route may
be more suitable when sequential covariates are observed; the choice
is about assumptions and target histories.

\section{Synthetic controls replace common trends by a factor restriction}
\label{r3:synthetic-control}
Suppose untreated outcomes obey
$Y_{it}(0)=\delta_t+\lambda_t' f_i+\varepsilon_{it}$, with
time-varying factors $\lambda_t$ and unit loadings $f_i$. A
time-invariant fixed effect is a special case. For one treated unit,
choose donor weights $w_j\ge0$, $\sum_jw_j=1$, using pre-intervention
outcomes and covariates only. The untreated prediction error is
\begin{equation}
 Y_{1t}(0)-\sum_jw_jY_{jt}(0)
 =\lambda_t'\!\left(f_1-\sum_jw_jf_j\right)
                 +\varepsilon_{1t}-\sum_jw_j\varepsilon_{jt}.
 \label{r3:synth-error}
\end{equation}
Exact loading balance removes the factor component at every time.
In a noiseless model, exact pretreatment fit identifies that balance
only when the preperiod factor matrix spans all relevant loading
directions. A factor emerging only after rollout defeats that inference.
This algebra explains the factor-model motivation and convex-hull
conditions in \citet[Section 2.2 and Appendix B]{abadie2010}.

For example, donors with loadings $0$ and $1$ cannot convexly match a
treated loading $2$. An unconstrained extrapolation can match it with
weights $(-1,2)$, at a different variance and robustness cost. A
penalized augmentation can reduce imbalance under an additional
outcome model, while an interactive-factor estimator estimates factors
and loadings directly. Those choices exchange interpolation stability
for model dependence; good pretreatment fit alone does not choose
between them. Donor contamination from a common promotion or balance
transfers also violates the required untreated donor interpretation.

With one treated region, placebo reassignment describes an inference
procedure only under an appropriate assignment or exchangeability
argument. It is not automatically an exact randomized-experiment
$p$-value. A historical fit across a handful of branches is likewise
not an independent sample of policy rollouts. The inferential regime
must be stated with the counterfactual model.

\section{Uncertainty follows the independent information, not the row count}
\label{r3:panel-inference}
For a differentiable estimating equation let $s_{it}$ be a scalar
linearized score. The variance of the sample mean is
\[
 \Var\!\left(\frac1{NT}\sum_{it}s_{it}\right)
 =\frac1{N^2T^2}\sum_{it,js}\Cov(s_{it},s_{js}).
\]
Clustering by unit retains within-unit terms but misses covariance
between branches facing common calendar shocks. Clustering by time
does the reverse. If dependence is permitted whenever either index
matches and absent otherwise, an inclusion--exclusion estimator adds
outer products of unit score sums and time score sums and subtracts
outer products of intersection-cell sums. This is a covariance model,
not a guarantee valid for arbitrary network or serial cross-unit
dependence.

For the exact random-effects score model
$s_{it}=a_i+b_t+\epsilon_{it}$ with mutually independent components,
\begin{equation}
 \Var(\bar s)=\frac{\sigma_a^2}{N}
             +\frac{\sigma_b^2}{T}
             +\frac{\sigma_\epsilon^2}{NT}.
 \label{r3:panel-information}
\end{equation}
Millions of branches cannot eliminate the finite-calendar uncertainty
represented by $\sigma_b^2/T$. If $b_t$ is serially dependent, replace
that term by $T^{-2}\sum_{t,s}\Cov(b_t,b_s)$. The period
clusters are no longer independent and a corresponding time-series
argument is needed. Nor does subtracting a common time effect remove
every common score component in a heterogeneous or nonlinear model.

Cross-fitting must respect the same structure. Holding out rows from
a customer leaves information from that customer's other periods in
training. Under two-way clustered sampling, a held-out rectangle can
train on observations outside both its row and column groups. The
multiway DML results of \citet[Assumptions 1--3, Theorems 1--2]{cluster2021}
require separately exchangeable, dissociated arrays, growing cluster
dimensions, score regularity and nuisance rates. Ordered banking time
series need not be separately exchangeable across calendar dates.
Those theorems cannot be invoked for serial dependence merely because
the software accepts two cluster identifiers.

\section{What the panel operation must declare}
A panel request must identify whether it targets a structural coefficient,
a cohort/time ATT, a path contrast or a donor-based counterfactual. It
must include the action's reversibility and carryover, the independent
unit, the time asymptotic regime, the treatment/control support and
the allowed common shocks. Fixed effects, lag instruments, trend
restrictions and factor balance are alternative pieces of scientific
evidence, not defaults inferred from the table's shape.

An eligible output preserves aggregation weights, unsupported horizons,
instrument restrictions, donor contamination checks and the covariance
construction. The bank can then distinguish an effect that changes
sign because of economics from a regression coefficient that changes
sign because it averages positive effects with negative weights.
That distinction is a prerequisite for financially meaningful
cross-application reuse.

\section{System moments, initial conditions and the practical panel choice}
\label{s5:dynamic-panel}
The difference equation removes an individual effect; it does not create a
strong instrument. This distinction matters for persistent loan balances or
capital ratios. Consider first the scalar model
\[
 y_{it}=\alpha_i+\rho y_{i,t-1}+v_{it},\qquad |\rho|<1,
 \quad \E[v_{it}\mid\alpha_i,y_{i1},v_{i2},\ldots,v_{i,t-1}]=0.
\]
Independent entities and fixed $T$ define the sampling experiment. For
$s\le t-2$, both $\E[y_{is}v_{it}]$ and $\E[y_{is}v_{i,t-1}]$ vanish, giving
$\E[y_{is}\Delta v_{it}]=0$. These are the differenced moments already
derived above \citep{arellano1991}. Their validity does not require the
initial observation to have the stationary distribution.

Level equations permit a different restriction. Write
$y_{i1}=\alpha_i/(1-\rho)+e_{i1}$. Then
\[
 \Delta y_{i2}=(\rho-1)e_{i1}+v_{i2},\qquad
 \E[\Delta y_{i2}\alpha_i]=0
 \quad\text{if }\E[e_{i1}\alpha_i]=0.
\]
The recursion
$\Delta y_{it}=\rho^{t-2}\Delta y_{i2}+
\sum_{j=0}^{t-3}\rho^j\Delta v_{i,t-j}$ propagates this zero covariance.
Consequently $\E[\Delta y_{i,t-1}(\alpha_i+v_{it})]=0$.
Combining these level moments with differenced moments yields system GMM.
Mean stationarity is sufficient for this construction; equality of initial
and limiting variances is not required. Conversely a portfolio assembled at
loan origination or after a balance-sheet repair may systematically start
above or below its entity-specific long-run level. The additional moment can
then fail even when the differenced moments remain valid. This reconstructs
the initial-condition argument in \citet[Sections 3.1--3.3, pp.107--109]{blundellbond2023}.
The archived text is their 2023 account, not the original 1998 article.

Adding $\beta x_{it}$ requires a separate classification for $x$.
Contemporaneous endogeneity excludes $x_{it}$ as an instrument; lagged levels
must precede both innovations in $\Delta v_{it}$. Predetermination and strict
exogeneity allow different lag sets. A level moment using $\Delta x_{i,t-1}$
also needs $\E[\Delta x_{i,t-1}\alpha_i]=0$ and orthogonality to $v_{it}$.
No initial-condition theorem for $y$ automatically supplies this condition
for a bank's chosen action. With common macro shocks, cross-entity
independence itself fails and the fixed-$T$, large-$N$ argument must change.

\subsection{Weighting, covariance and what a correction can repair}
Let $b_i-G_i\theta$ collect the entity's moments after any prespecified lag
restriction and compression, where $\theta=(\rho,\beta)'$. For a positive
weight $W$, the quadratic criterion has minimizer
\[
 \hat\theta_W=(\bar G'W\bar G)^{-1}\bar G'W\bar b,
 \quad
 \sqrt N(\hat\theta_W-\theta_0)=
 (G'WG)^{-1}G'W\,N^{-1/2}\sum_i(b_i-G_i\theta_0)+o_p(1).
\]
The expansion needs $\bar G\to G$, a moment CLT, a nonsingular $G'WG$,
and consistent $W$. Thus the covariance estimator is
\[
 \widehat{\Var}(\hat\theta_W)=N^{-1}
 (\bar G'W\bar G)^{-1}\bar G'W\hat\Omega W\bar G
 (\bar G'W\bar G)^{-1},
\]
where $\hat\Omega$ is the covariance of the entire entity moment vector,
not of individual entity-date rows. Two-step GMM takes
$W=\hat\Omega(\hat\theta_1)^{-1}$ using an initial consistent estimator.
The first-order limit omits uncertainty in that estimated weight because
the population moment is zero. The finite-sample derivative is nevertheless
nonzero:
\[
 d\hat\theta_2=(\bar G'W\bar G)^{-1}\bar G'(dW)
 (\bar b-\bar G\hat\theta_2),\qquad
 dW=-W(d\hat\Omega)W.
\]
This displays the term whose variation a finite-sample correction must
address. The 2023 primary account explains the Windmeijer correction and its
motivation \citep[Section 2.3, p.105]{blundellbond2023}. The transparent local
reference implements the first-order sandwich, \emph{not} that correction;
the correction's original paper has not been substituted by a metadata record.
A production two-step adapter needs a source-qualified implementation of the
correction or a separately justified complete-refit alternative. Neither can
make invalid moments valid or rescue a rank-deficient Jacobian.

Compression also changes a statistical problem. An unrestricted triangular
lag matrix has $(T-2)(T-1)/2$ instruments in the scalar differenced model.
The reference collapses moments over dates and keeps lags two and three,
yielding four differenced moments for $(y,x)$, or six after adding level
moments. Scaling moments by observed root mean squares changes the one-step
weight; it leaves the population zero restrictions unchanged. The code
therefore reports this estimator as scaled identity-weight GMM, not the
canonical homoskedastic optimal one-step estimator. Small objective values
with many instruments can reflect overfitting. An overidentification test
is neither proof of exclusion nor evidence of adequate instrument strength.

\subsection{A paired comparison and a banking decision rule}
The synthetic short-panel experiment uses $N=800$, eight transitions after
an initial observation, and
$y_{it}=\alpha_i+\rho y_{i,t-1}+.6x_{it}+v_{it}$, with
$x_{it}=.5x_{i,t-1}+.4v_{it}+\eta_{it}$. A 400-period burn-in generates
the regular initial law. The initial-condition challenger then adds a
multiple of $\alpha_i$ to the initial balance; the serial-error challenger
changes the innovation law. The four methods use identical datasets.
Coverage below is joint for both coefficients, using two-coordinate
Bonferroni 95\% intervals across 400 repetitions.
\begin{center}\small
\begin{tabular}{lrrrr}\toprule
Regime & Difference 1 & Difference 2 & System 1 & System 2\\\midrule
Stationary & .955 & .950 & .958 & .948\\
Persistent & .985 & .983 & .968 & .933\\
Initial departure & .950 & .960 & .033 & .000\\
Serial-error violation & .000 & .000 & .000 & .000\\\bottomrule
\end{tabular}\end{center}
One and two denote weighting steps, not different causal targets. Near
95\% coverage, the Monte Carlo standard error is about .011. The very high
coverage of a weak differenced design is not a precision advantage; weak
moments can produce broad intervals and unreliable normal approximations.
The initial-departure failure is the decisive exclusion from an automatic
system-GMM default. For established comparable entities, system moments can
be a challenger after a substantive initial-law argument. For new portfolios,
the primary differenced route has the weaker initial-law requirement, but
must still report Jacobian strength, instrument count and residual serial
correlation. If those fail, the result is an unsupported inferential claim,
not a prompt to search until a small standard error appears.

Misspecification-robust moment inference is another distinct alternative.
\citet[Sections 2--3, Assumptions 1--2]{kleibergen2021} targets the minimizer
of a continuously updated population GMM criterion when no parameter makes
all moments zero. Its joint moment/Jacobian CLT and covariance conditions
support an identification- and misspecification-robust score construction.
That pseudo-true minimizer need not equal a causal response or a meaningful
risk premium. This use of ``double robust'' concerns weak identification and
moment misspecification; it is different from outcome/propensity robustness
in Chapter~\ref{r3:learning}. The inspected copy is the 2021 preprint of work
identified in 2025 journal metadata. It is a substantive inference challenger,
not a license to keep a causal label after the identifying restrictions fail.


\section{Normalized cohort effects, aggregation and a selected donor test}
\label{s6:panel-influence}
The normalized representation in the earlier panel section has uncertainty
from both its numerators and its denominators. For a fixed cohort/date pair,
write $G$ for membership of the cohort, $C$ for eligible controls,
$D=Y_t-Y_{g-1}$, $p(x)=P(G=1\mid G+C=1,X=x)$ and
$m(x)=\E[D\mid C=1,X=x]$. Require disjoint $G,C$, no anticipation,
conditional parallel trends, overlap in the eligible population, and the
same cohort/calendar target as in \citet{callaway2021}. Other cohorts may
have $G=C=0$. Put $b=Cp/(1-p)$, $p_G=\E G$, $p_b=\E b$ and
\[
 \mu_G=\frac{\E[G(D-m)]}{p_G},\qquad
 \mu_b=\frac{\E[b(D-m)]}{p_b},\qquad \tau=\mu_G-\mu_b.
\]
At the true nuisances, $p_b=p_G$, $\mu_b=0$ and $\mu_G=\tau$.
The sample estimator normalizes each weighted mean separately. Holding
nuisances fixed, quotient differentiation yields
\begin{equation}
 \varphi_{gt}=
 \frac{G}{p_G}(D-m-\mu_G)
 -\frac{b}{p_b}(D-m-\mu_b).
 \label{s6:att-if}
\end{equation}
In particular, the treated term contains $-G\tau/p_G$. Centering an
unnormalized score after the fact does not reproduce this random-cohort
contribution.

At the true nuisances, the derivative in $m$ is
$-\E[G\,dm]/p_G+\E[b\,dm]/p_G=0$. The propensity derivative
multiplies a control residual of conditional mean zero, including the
normalizer derivative. These equalities justify the orthogonal influence
under cross-fitting and the appropriate product-rate conditions. With just
one nuisance correct, point consistency can survive, but those two
derivatives need not both vanish at the fitted limits. A doubly robust point
estimate is not automatically a doubly robust standard error; either obtain
the stronger rates/correct limits or include the first-stage influence of
the misspecified-limit procedure. Cluster the full influence by the actual
independent sampling unit. Many customers observed under one branch rollout
do not create independent branch assignments.

For an event-time average over a fixed eligible cohort set $\mathcal G$,
let $P=\sum_{g\in\mathcal G}p_g$ and
$\delta=\sum_gp_g\tau_{g,g+e}/P$. Its influence is
\begin{equation}
 \varphi_\delta=
 \sum_g\frac{p_g}{P}\varphi_{g,g+e}
 +\frac1P\sum_g G_g(\tau_{g,g+e}-\delta).
 \label{s6:aggregate-att-if}
\end{equation}
The second term accounts for estimated cohort shares. It disappears only
when those shares are fixed externally or the component effects coincide.
The first term is one sum on the same observation: controls reused across
cohorts create covariance. Forming a diagonal matrix of separately reported
standard errors loses that covariance. If eligibility itself is selected
using estimated effects, this fixed-set derivation is no longer applicable.
Financial exposure weighting adds the analogous derivative of estimated
exposures; externally fixed accounting weights use the full covariance of
the component effects without that extra term.

\subsection{What an exact synthetic-control inference route actually needs}
Synthetic-control weights interpolate pre-intervention outcomes, usually
under nonnegative weights summing to one. Their value depends on pre-fit,
factor balance, donor contamination and whether the treated unit is inside
the donor support \citep{abadie2010,abadie2021,benmichael2021,synthdid2021}.
A fitted counterfactual curve is not a confidence procedure by itself.
Consider one precise banking design: one of $J$ eligible regions is selected
uniformly for a sustained pilot, after the eligibility set and fitting
algorithm are fixed. Under the sharp null of no effect in every post-period,
the complete untreated outcome matrix is invariant to this assignment.

For each candidate treated region $j$, refit its donor weights using its
own pre-period and the remaining $J-1$ regions. Apply the same fixed
regularization and fit exclusions to every candidate. Compute a discrepancy
$T_j$, such as post-period mean squared gap divided by a pre-period error
with a fixed positive floor. With observed assignment $a$,
\[
 p_a=J^{-1}\sum_{j=1}^J\mathbf1\{T_j\ge T_a\}
\]
is super-uniform under the randomization: conditional on the fixed outcome
matrix, $a$ is uniform over the $J$ positions in its rank list. Ties make
the displayed test conservative. Serial correlation, common shocks and
weight estimation do not invalidate this finite assignment argument because
the whole paths are held fixed and the fitting algorithm is applied
symmetrically. Few eligible regions impose coarse attainable significance.

For a sharp constant additive effect $\tau$ in every post-period, subtract
$\tau$ from the actual treated path to impute the untreated matrix. For
each possible assignment add the hypothesized effect to that region, then
recompute the discrepancy centered on $\tau$, including its donor fit.
Inverting the resulting tests gives a confidence set for this \emph{sharp
constant-effect model}. It is not automatically an interval for an arbitrary
average treatment effect under heterogeneous post-period effects.

An observationally selected distressed bank is not uniformly randomized
among healthy donors. Its donor placebo ranks are then descriptive unless
another exchangeability/design argument is supplied. A residual-permutation
or Bayesian structural time-series route substitutes its own temporal or
prior/model restrictions; neither inherits the assignment proof. This is
why the product selects the above test for randomized regional pilots and
returns an explicit additional-design requirement for observational donor
studies. The method choice concerns the source of uncertainty, not a claim
that one donor-fitting algorithm is universally preferable.


A stronger temporal challenger is the null-imposed conformal construction
of \citet{conformal2022}. It imputes the hypothesized untreated outcome,
refits the counterfactual prediction under that null, and ranks its residual
discrepancy over a declared permutation group. Exact validity follows under
the required exchangeability and permutation-equivariant fitting. For
dependent observations, their Theorem 1 instead assumes a fixed postperiod
length, stationary mixing shocks with a bounded discrepancy density and
small prediction errors both on average and at the tested dates. Theorem 2
uses estimator stability plus a stationary exponentially mixing data law.
These are substantive ways to trade assignment information for temporal
prediction restrictions. Neither is a license to permute arbitrary
nonstationary bank residuals, and neither automatically changes a sharp
path hypothesis into the mean effect in Equation~\eqref{s6:aggregate-att-if}.
For an observational regional rollout this challenger can be selected only
after its null prediction and temporal restrictions are defended. The
randomized pilot route was selected here because its source of repeated
assignment is specified directly.

